\ifx\allfiles\undefined
\documentclass[12pt, a4paper, oneside, UTF8]{ctexbook}
\def\path{../config}
\input{\path/package.tex}

% 定理环境设置
\input{\path/theorem1_zh.tex}

\input{\path/custom.tex}

% 修改参数改变封面样式，0 默认原始封面、内置其他1、2、3种封面样式
\def\myIndex{3}


\ifnum\myIndex>0
    \input{\path/cover_package_\myIndex} 
\fi

\def\myTitle{线性代数笔记}
\def\myAuthor{M.K.}
\def\myDateCover{\today}
\def\myDateForeword{\today}
\def\myForeword{前言}
\def\myForewordText{
    
    这是一个基于\LaTeX{}模板的线代笔记。

}
\def\mySubheading{副标题}


\begin{document}
% \input{\path/cover_text_\myIndex.tex}

\newpage
\thispagestyle{empty}
\begin{center}
    \Huge\textbf{\myForeword}
\end{center}
\myForewordText
\begin{flushright}
    \begin{tabular}{c}
        \myDateForeword
    \end{tabular}
\end{flushright}

\newpage
\pagestyle{plain}
\setcounter{page}{1}
\pagenumbering{Roman}
\pagestyle{fancy}
\fancyfoot[C]{\thepage}
\renewcommand{\headrulewidth}{0.4pt}
\renewcommand{\footrulewidth}{0pt}
\tableofcontents

\newpage
\pagenumbering{arabic}
\setcounter{page}{1}








\else
\fi
\chapter{线性方程组}

\section{线性方程组有解的充要条件}
\begin{theorem}
    设$\bm{A}$为$m\times n$矩阵，$\bm{b}\in \mathbb{R}^m$，则线性方程组$\bm{Ax}=\bm{b}$有解的充分必要条件是$\rank(\bm{A}) = \rank(\bm{A}|\bm{b})$．
\end{theorem}
\section{线性方程组的解的结构}
\begin{theorem}[齐次线性方程组解的结构]
    设$\bm{A}$为$m\times n$矩阵，$\rank(\bm{A}) = r$，则齐次线性方程组$\bm{Ax}=\bm{0}$的解集构成$\R^n$的一个子空间（称为\textbf{解空间}），其维数为$n-r$．
    若$\bm{\xi}_1,\bm{\xi}_2,\ldots,\bm{\xi}_{n-r}$为解空间的一组基（即\textbf{基础解系}），则通解为
    \[
        \bm{x} = k_1\bm{\xi}_1 + k_2\bm{\xi}_2 + \cdots + k_{n-r}\bm{\xi}_{n-r},\quad k_1,k_2,\ldots,k_{n-r}\in\R
    \]
\end{theorem}

\begin{theorem}[非齐次线性方程组解的结构]
    设$\bm{A}$为$m\times n$矩阵，$\rank(\bm{A}) = \rank(\bm{A}|\bm{b}) = r$，即$\bm{Ax}=\bm{b}$有解．
    若$\bm{x}_0$为$\bm{Ax}=\bm{b}$的一个特解，$\bm{\xi}_1,\bm{\xi}_2,\ldots,\bm{\xi}_{n-r}$为对应齐次方程组$\bm{Ax}=\bm{0}$的基础解系，则通解为
    \[
        \bm{x} = \bm{x}_0 + k_1\bm{\xi}_1 + k_2\bm{\xi}_2 + \cdots + k_{n-r}\bm{\xi}_{n-r},\quad k_1,k_2,\ldots,k_{n-r}\in\R
    \]
\end{theorem}

\begin{example}
    设$\bm{A}_{4\times 4}=(\bm{\alpha_1} , \bm{\alpha_2} , \bm{\alpha_3} , \bm{\alpha_4})$，
    且$\bm{Ax}=\bm{0}$的通解为$k(1 , 0 , -4 , 0)^T$，则下列是$\bm{A}^*\bm{x}=\bm{0}$的基础解系的是（~~~~~~）
    \begin{tasks}[style = itemize,label = \Alph*.](4)
    \task $\alpha_1,\alpha_2,\alpha_3$
    \task $\alpha_1,\alpha_2$
    \task $\alpha_1,\alpha_2,\alpha_4$
    \task $\alpha_1,\alpha_3,\alpha_4$
    \end{tasks}
\end{example}
\begin{solution}
    由通解可知，$\rank(\bm{A}) = 3$，$\bm{A}^*\bm{A}=|\bm{A}|\bm{E}=\bm{0}\Rightarrow\bm{A}^*\bm{\alpha_1}=\bm{A}^*\bm{\alpha_2}=
    \bm{A}^*\bm{\alpha_3}=\bm{A}^*\bm{\alpha_4}=\bm{0}$，并且$\rank(\bm{A}^*)=1$，故$\bm{A}^*\bm{x}=\bm{0}$的基础解系有3个解向量，
    又$\bm{\alpha_1}-4\bm{\alpha_3}=\bm{0}\Rightarrow\bm{\alpha_1}\text{与}\bm{\alpha_3}$线性相关，故选C
\end{solution}

\section{解方程组的基本方法}
\begin{theorem}[基础解系的求法]
    设 $\bm{A}$ 为 $m\times n$ 矩阵，且 $\rank(\bm{A})=r$，则齐次线性方程组
    \[
    \bm{Ax}=\bm{0}
    \]
    的解空间维数为 $n-r$，其基础解系由 $n-r$ 个线性无关的解向量组成。
    \textbf{求解步骤：}
    \begin{enumerate}[label = (\arabic*)]
        \item \textbf{消元化简：}  
        对 $\bm{A}$ 作初等行变换，将其化为行最简形（RREF），以便直接读出变量关系。

        \item \textbf{确定变量类型：}  
        在行最简形中，每一行的主元所在列对应\textbf{主变量}；其余变量为\textbf{自由变量}，自由变量的个数为 $n-r$。

        \item \textbf{引入参数：}  
        用自由变量表示主变量，将解写成参数形式（解空间的通解表达）。

        \item \textbf{构造基础解向量：}  
        令每个自由变量依次取 $1$，其余取 $0$（即标准基），代入通解，得到一组解向量。

        \item \textbf{得到基础解系：}  
        所有得到的解向量线性无关，共 $n-r$ 个，即构成基础解系。
    \end{enumerate}
\end{theorem}

\begin{example}
    设$\bm{A}_{4\times 4}=(\bm{\alpha_1} , \bm{\alpha_2} , \bm{\alpha_3} , \bm{\alpha_4})$，
    且$\bm{Ax}=\bm{\beta}$的通解为$k(1 , -2 , 4 , 0)^T+(1, 2, 2,1)^T$，记$\bm{B}=(\bm{\alpha_3}, \bm{\alpha_2}, \bm{\alpha_1}, 
    \bm{\beta}-\bm{\alpha_4})$，求$\bm{Bx}=\bm{\alpha_1}-\bm{\alpha_2}$的通解．
\end{example}
\begin{solution}
    由题意得到$\rank(\bm{A})=3$，$\bm{\alpha_1}-2\bm{\alpha_2}+4\bm{\alpha_3}=\bm{0},\bm{\alpha_1}+\bm{\alpha_2}+2\bm{\alpha_3}+\bm{\alpha_4}=\bm{\beta}
    \Rightarrow (4, -2, 1, 0)^T,(2, 2, 1, -1)^T$是$\bm{Bx}=\bm{0}$的解，且二者线性无关．
    显然$(0, -1, 1, 0)^T$是$\bm{Bx}=\bm{\alpha_1}-\bm{\alpha_2}$的一个特解，$\rank(\bm{B})=2$，于是可知$\bm{Bx}=\bm{\alpha_1}-\bm{\alpha_2}$的通解为
    \[
        \bm{x} = k_1(4, -2, 1, 0)^T + k_2(2, 2, 1, -1)^T + (0, -1, 1, 0)^T
    \]
\end{solution}

\section{练习}
\begin{exercise}
    设矩阵$A=\begin{bNiceMatrix}
    0 & -1 & 2\\
    -1 & 0 & 2\\
    -1 & -1 & 3
    \end{bNiceMatrix}$
    ，求所有满足
    $\begin{cases}
        A^2\alpha=\alpha+2\beta\\
        A\alpha=\alpha+\beta
    \end{cases}$
    的向量 $\alpha,\beta$
\end{exercise}
\begin{solution}
    将 $\beta = (A - E)\alpha$ 代入 $A^2\alpha = \alpha + 2\beta$得
    \[A^2\alpha = \alpha + 2(A - E)\alpha\implies (A - E)^2\alpha = \bm{0}\]
    计算
    \[
        A - E = \begin{bNiceMatrix}
            -1 & -1 & 2\\
            -1 & -1 & 2\\
            -1 & -1 & 2
        \end{bNiceMatrix},\qquad
        (A - E)^2 = \begin{bNiceMatrix}
            -1 & -1 & 2\\
            -1 & -1 & 2\\
            -1 & -1 & 2
        \end{bNiceMatrix}^2 = \bm{0}
    \]
    故 $(A - E)^2\alpha = \bm{0}$ 对任意 $\alpha\in\R^3$ 成立，进而 $\beta = (A - E)\alpha$，即
    \[\boxed{\alpha = \begin{bNiceMatrix}
        x_1\\
        x_2\\
        x_3
    \end{bNiceMatrix},\quad \beta = \begin{bNiceMatrix}
        -x_1 - x_2 + 2x_3\\
        -x_1 - x_2 + 2x_3\\
        -x_1 - x_2 + 2x_3
    \end{bNiceMatrix},\quad x_1,x_2,x_3\in\R}\]
\end{solution}

\newcounter{inleq}
\setcounter{inleq}{1}
\begin{exercise}
    已知方程组(\Roman{inleq})
    $\begin{cases}
        x_1+2x_2+3x_3=0\\
        2x_1+3x_2+5x_3=0\\
        x_1+x_2+ax_3=0
    \end{cases}$
\refstepcounter{inleq}
    与(\Roman{inleq})
    $\begin{cases}
        x_1+bx_2+cx_3=0\\
        2x_1+b^2x_2+(c+1)x_3=0\\
    \end{cases}$
    有相同的解集，求$a,b,c$的值．
\end{exercise}
\setcounter{inleq}{1}
\begin{solution}
    记方程组(\Roman{inleq})\refstepcounter{inleq}(\Roman{inleq})的系数矩阵分别为$\bm{A},\bm{B}$，则
    \[\rank(\bm{A})=\rank(\bm{B})
    =\rank\begin{bNiceMatrix}
        \bm{A}\\
        \bm{B}
    \end{bNiceMatrix}\text{且}\rank(\bm{B})\leq 2\Rightarrow\rank(\bm{A})\leq 2\]
    故$|\bm{A}|=0\Rightarrow a=2\Rightarrow\rank(\bm{A})=2\Rightarrow\rank(\bm{B})
    =\rank\begin{bNiceMatrix}
        \bm{A}\\
        \bm{B}
    \end{bNiceMatrix}=2$，于是
    \[
        \begin{bNiceMatrix}
            \bm{A}\\
            \bm{B}
        \end{bNiceMatrix}=\begin{bNiceMatrix}
            1 & 2 & 3\\
            2 & 3 & 5\\
            1 & 1 & 2\\
            1 & b & c\\
            2 & b^2 & c+1
        \end{bNiceMatrix}\rightarrow\begin{bNiceMatrix}
            1 & 2 & 3\\
            0 & 1 & 1\\
            0 & 0 & c-b-1\\
            0 & 0 & -bc-c+3b+1\\
            0 & 0 & 0
        \end{bNiceMatrix}\Rightarrow\begin{cases}
            c-b-1=0\\
            -bc-c+3b+1=0
        \end{cases}
    \]
    解得$\begin{cases}
        b=0\\
        c=1
    \end{cases}$或$\begin{cases}
        b=1\\
        c=2
    \end{cases}$，经检验，$b=0,c=1$时，$\rank({\bm{B}})=1$，故舍去，最终结果为\[a=2,b=1,c=2\]
\end{solution}
\ifx\allfiles\undefined
\end{document}
\fi